\documentclass[11pt]{article}
\usepackage[margin=1in]{geometry}
\usepackage{amsmath,amssymb,amsthm,booktabs,array,longtable}
\usepackage[T1]{fontenc}
\usepackage{lmodern}
\usepackage[hidelinks]{hyperref}
\newtheorem{theorem}{Theorem}
\newtheorem{corollary}[theorem]{Corollary}
\newcommand{\N}{\mathbb N}
\title{Weighted extensions of Jovovic's 2003 formula:\protect\\
resolving Mathar's conjecture\protect\\
(OEIS A084261, FORMULA, 30 November 2012)}
\author{Haobo Ma\\\small ChronoAI Pte Ltd, Singapore
\and Wenlin Zhang\\\small National University of Singapore\\
\small ChronoAI Pte Ltd, Singapore}
\date{8 November 2026}
\begin{document}
\maketitle
\begin{abstract}
We give an alternative proof of Jovovic's 2003 exponential generating
function and resolve Mathar's ``Conjecture'' in OEIS A084261,
FORMULA, 30 November 2012, for every integer $n\geq3$.
For $a_n(t)=\sum_{k=0}^{\lfloor n/2\rfloor}\binom n{2k}k!t^{n-2k}$,
the fixed-point-weighted generating function of Fang, Hamaker and
Troyka (2020; journal version 2022) yields an explicit polynomial recurrence.
A common moment representation also proves the recorded row sums,
continued fraction, Hankel product and asymptotic equivalent, the last
extending to every fixed real $t>0$.
\end{abstract}

\section{The conjecture and its proof}
The OEIS entry A084261, introduced by Paul Barry in 2003, records Mathar's
conjecture as a relation between four successive terms of the binomial
transform of factorial numbers with interpolated zeros~\cite{oeis}.
The following theorem supplies its full nonnegative-index scope.
The later enumeration theorem of Fang, Hamaker and Troyka
\cite[Proposition 4.13 and following example]{fang} supplies the weighted
generating function used in Section~2, but does not state this recurrence.
The source gives no starting index; $n\geq3$ makes every index in the displayed
recurrence nonnegative.

\begin{theorem}\label{thm:main}
For $n\in\N=\{0,1,2,\ldots\}$ put
\begin{equation}\label{eq:sum}
 a_n=\sum_{k=0}^{\lfloor n/2\rfloor}\binom n{2k}k!.
\end{equation}
Then $a_0=a_1=1$, $a_2=2$, and, for every integer $n\geq3$,
\begin{equation}\label{eq:rec}
 2a_n-4a_{n-1}+(2-n)a_{n-2}+(n-2)a_{n-3}=0.
\end{equation}
\end{theorem}

\begin{proof}
Define $b_{2j}=j!$, $b_{2j+1}=0$ for $j\geq0$, and, for any sequence $c$, set
\[
 T(c,m)=\sum_{k=0}^{m}\binom mk c_k,\qquad
 T_j(m)=T((b_{k+j})_{k\geq0},m).
\]
Thus $T_0(m)=a_m$. Pascal's identity, including its zero endpoint terms, gives
\begin{equation}\label{eq:pascal}
 T_j(m+1)=T_j(m)+T_{j+1}(m)\qquad(m,j\geq0).
\end{equation}
The seed sequence satisfies
\begin{equation}\label{eq:seed}
 2b_{k+3}+2b_{k+2}-3b_{k+1}-2b_k
       =k(b_{k+1}+b_k)\qquad(k\geq0).
\end{equation}
Indeed, for $k=2j$ the two sides are respectively
$2(j+1)!-2j!=2j\,j!$ and $2j\,j!$.
For $k=2j+1$ they are
$2(j+2)!-3(j+1)!=(2j+1)(j+1)!$ and $(2j+1)(j+1)!$.

We next use the weighted identity
\begin{equation}\label{eq:weighted}
 mT_1(m)=\sum_{k=0}^{m}\binom mk k b_{k+1}
          +\sum_{k=0}^{m}\binom mk k b_k.
\end{equation}
To see this, use
$(m-k)\binom mk=(k+1)\binom m{k+1}$, taking
$\binom m{m+1}=0$, and reindex:
\[
 \sum_{k=0}^{m}(m-k)\binom mk b_{k+1}
 =\sum_{j=1}^{m+1}j\binom mj b_j
 =\sum_{j=0}^{m}j\binom mj b_j.
\]
Adding $\sum_{k=0}^{m}k\binom mk b_{k+1}$ proves
\eqref{eq:weighted}, also when $m=0$.
Multiply \eqref{eq:seed} by $\binom mk$, sum, and use
\eqref{eq:weighted}. The result is
\[
 2T_3(m)+2T_2(m)-(m+3)T_1(m)-2T_0(m)=0.
\]
By \eqref{eq:pascal},
\begin{align*}
 T_1(m)&=a_{m+1}-a_m,\\
 T_2(m)&=a_{m+2}-2a_{m+1}+a_m,\\
 T_3(m)&=a_{m+3}-3a_{m+2}+3a_{m+1}-a_m.
\end{align*}
Substitution gives
$2a_{m+3}-4a_{m+2}-(m+1)a_{m+1}+(m+1)a_m=0$.
Set $m=n-3$. The initial values follow directly from \eqref{eq:sum}.
\end{proof}

The recurrence and the three initial values determine the sequence uniquely:
the coefficient of $a_n$ is the nonzero constant $2$. At the first two
admissible indices it gives $a_3=4$ and $a_4=9$.

\section{Generating function and weighted recurrence}
For a scalar $t$ define the polynomial
\begin{equation}\label{eq:polynomial}
 a_n(t)=\sum_{k=0}^{\lfloor n/2\rfloor}\binom n{2k}k!\,t^{n-2k}.
\end{equation}
Its exponential generating function is
\begin{equation}\label{eq:egf}
 F_t(z)=\sum_{n\geq0}a_n(t)\frac{z^n}{n!}
 =e^{tz}B(z),\qquad
 B(z)=\sum_{k\geq0}\frac{k!z^{2k}}{(2k)!}.
\end{equation}
This follows by multiplying the two series and comparing coefficients.
Fang, Hamaker and Troyka~\cite[Proposition 4.13 and its following example,
arXiv v3, p.~16]{fang} prove this factorization with $t$ marking fixed points in
their enumeration of $2143$-avoiding involutions. Their interpretation uses
their involution-containment relation, rather than ordinary permutation
pattern containment. The proof and example already appear as
Theorem~1.7 and Section~4.3 in their 31 August 2020 first version.
The next corollary extracts a differential recurrence from that known
factorization and reproves Jovovic's recorded formula.

\begin{corollary}\label{cor:weighted}
For every $n\geq3$, as an identity of polynomials in $t$,
\begin{equation}\label{eq:weightedrec}
 2a_n(t)-4t\,a_{n-1}(t)+(2t^2-n)a_{n-2}(t)
       +t(n-2)a_{n-3}(t)=0.
\end{equation}
The initial polynomials are $1,t,t^2+1$.
Moreover, for every complex $z$,
\begin{equation}\label{eq:erf}
 F_1(z)=e^z\left(1+\frac{\sqrt\pi}{2}z e^{z^2/4}
                                  \operatorname{erf}(z/2)\right),
\end{equation}
where $\operatorname{erf}(w)=\frac{2}{\sqrt\pi}\int_0^w e^{-u^2}\,du$.
\end{corollary}

\begin{proof}
The ratio of successive absolute terms in $B(z)$ is
$|z|^2(k+1)/((2k+2)(2k+1))$, which tends to zero. Thus $B$ is entire.
Comparing its coefficients gives
\begin{equation}\label{eq:Bode}
 2B''(z)-zB'(z)-2B(z)=0,\qquad B(0)=1,\quad B'(0)=0.
\end{equation}
For the coefficient of $z^{2k}$, the first term is
$2(k+1)!/(2k)!$, while the other two together are
$(2k+2)k!/(2k)!$; all odd coefficients vanish.
Substitute $B=e^{-tz}F_t$ to obtain
\[
 2F_t''-(4t+z)F_t'+(2t^2+tz-2)F_t=0.
\]
For $m\geq1$ its coefficient of $z^m/m!$ is
\[
 2a_{m+2}(t)-4t\,a_{m+1}(t)+(2t^2-2-m)a_m(t)
       +tm\,a_{m-1}(t)=0.
\]
Set $m=n-2$ to get \eqref{eq:weightedrec}. The constant coefficient
also gives $2a_2(t)-4t\,a_1(t)+(2t^2-2)a_0(t)=0$.

For \eqref{eq:erf}, set
$G(z)=e^{z^2/4}\int_0^z e^{-u^2/4}\,du$ and $C(z)=1+zG(z)/2$.
The identity $G'=zG/2+1$ yields
\[
 C'=(1/2+z^2/4)G+z/2,\qquad
 C''=(3z/4+z^3/8)G+1+z^2/4.
\]
Consequently $2C''-zC'-2C=0$, $C(0)=1$, and $C'(0)=0$.
The differential equation determines successive Taylor coefficients
uniquely, so $C=B$. Finally,
$\int_0^z e^{-u^2/4}\,du=\sqrt\pi\operatorname{erf}(z/2)$.
Equation \eqref{eq:erf} recovers Jovovic's formula in the OEIS entry,
dated 25 September 2003~\cite{oeis}.
\end{proof}

\section{Moment consequences: row sums, Hankel determinants and fractions}
The same weighted family has a positive moment representation:
\begin{equation}\label{eq:moment}
 a_n(t)=\int_{\mathbb R}(t+x)^n|x|e^{-x^2}\,dx
 \qquad(n\geq0,\ t\in\mathbb R).
\end{equation}
Indeed the measure has total mass one and its odd moments vanish.
Its $2k$th moment is
\[
 \int_0^\infty u^ke^{-u}\,du=k!.
\]
Expansion of $(t+x)^n$ proves \eqref{eq:moment}.

\begin{corollary}\label{cor:triangle}
The row polynomial of OEIS A161556 is $a_n(t)$; in particular its row
sum is $a_n(1)$. For the array $T(n,j)$, with $0\leq j\leq n$,
\[
 T(n,j)=\begin{cases}\binom nj((n-j)/2)!,&n-j\text{ even},\\
 0,&n-j\text{ odd},\end{cases}\qquad
 \sum_{j=0}^n T(n,j)t^j=a_n(t).
\]
Its column exponential generating function is $z^jB(z)/j!$.
\end{corollary}
\begin{proof}
Substitute $j=n-2k$ in the row sum. For the column, put $n=j+2k$:
\[
 \binom{j+2k}{j}\frac{k!z^{j+2k}}{(j+2k)!}
 =\frac{z^j}{j!}\frac{k!z^{2k}}{(2k)!}.
\]
This also gives the exponential Riordan description $[B(z),z]$
in A161556~\cite{oeis-triangle}.
\end{proof}

\begin{corollary}\label{cor:hankel}
For every real $t$ and $n\geq0$,
\begin{equation}\label{eq:hankel}
 H_n(t):=\det(a_{i+j}(t))_{0\leq i,j\leq n}
 =\prod_{r=1}^n\left\lceil\frac r2\right\rceil^{n+1-r}>0.
\end{equation}
The ordinary generating series, interpreted formally, is
\begin{equation}\label{eq:cf}
 \sum_{n\geq0}a_n(t)z^n
 =\cfrac{1}{1-tz-\cfrac{\beta_1z^2}{1-tz-
       \cfrac{\beta_2z^2}{1-tz-\cfrac{\beta_3z^2}{1-tz-\cdots}}}},
 \qquad \beta_r=\left\lceil\frac r2\right\rceil.
\end{equation}
At $t=1$, \eqref{eq:cf} is Hanna's 17 January 2006 fraction in A084261;
\eqref{eq:hankel} is the product in A137704. The bivariate version
of \eqref{eq:cf} is the fraction in A161556.
\end{corollary}
\begin{proof}
Write $L_m^{(\alpha)}$ for the Laguerre polynomial with Rodrigues formula
\[
 L_m^{(\alpha)}(u)=\frac{u^{-\alpha}e^u}{m!}
       \frac{d^m}{du^m}(e^{-u}u^{m+\alpha}),\qquad\alpha=0,1.
\]
Its leading coefficient is $(-1)^m/m!$. For $j<m$, integration by
parts $m$ times gives
\[
 \int_0^\infty u^{\alpha+j}e^{-u}L_m^{(\alpha)}(u)\,du=0.
\]
All boundary terms vanish since the differentiated factor has a zero
of sufficient order at zero and exponential decay at infinity.
For $j=m$ the same calculation gives
$(-1)^m\Gamma(m+\alpha+1)$. Multiplying by the leading coefficient
therefore gives
\[
 \int_0^\infty u^\alpha e^{-u}(L_m^{(\alpha)}(u))^2\,du
 =\frac{\Gamma(m+\alpha+1)}{m!}.
\]
These normalizations agree with the standard Laguerre
conventions~\cite[Sections 18.3 and 18.5]{dlmf}.
Thus the monic polynomials orthogonal for $|x|e^{-x^2}\,dx$ are
\[
 P_{2m}(x)=(-1)^mm!L_m^{(0)}(x^2),\qquad
 P_{2m+1}(x)=(-1)^mm!xL_m^{(1)}(x^2).
\]
Their squared norms are respectively
$h_{2m}=(m!)^2$ and $h_{2m+1}=m!(m+1)!$.
For the measure translated by $t$, use $Q_r(y)=P_r(y-t)$.
The monic change of basis from $1,y,\ldots,y^n$ to $Q_0,\ldots,Q_n$
is triangular with determinant one. Consequently its moment determinant
is $\prod_{r=0}^nh_r$. Since $h_0=1$ and
$h_r/h_{r-1}=\lceil r/2\rceil$, this is \eqref{eq:hankel}.
It equals $\prod_{k=0}^n\lfloor(k+2)/2\rfloor^{n-k}$, the product
recorded in A137704~\cite{oeis-hankel}.

Orthogonality and parity give the three-term relation
\[
 yQ_r=Q_{r+1}+tQ_r+\beta_rQ_{r-1}\quad(r\geq0),
\]
where $Q_{-1}=0$ and $\beta_0=0$.
Indeed all coefficients below $Q_{r-1}$ vanish by orthogonality;
the diagonal coefficient before translation is zero by parity, and
the coefficient of $Q_{r-1}$ is $h_r/h_{r-1}$ by symmetry of
multiplication by $y$ in the inner product.
The coefficient of $Q_0$ in $y^nQ_0$ is $a_n(t)$.
Expanding the relation repeatedly expresses it as a sum over
nonnegative paths of length $n$ from height zero to zero: an up step
has weight one, a level step weight $t$, and a down step from height
$r$ weight $\beta_r$. If $M_r(z)$ counts paths at or above height $r$
that start and finish at $r$, decomposition at their first return gives
\[
 M_r=1+tzM_r+\beta_{r+1}z^2M_{r+1}M_r.
\]
Solving successively for $M_r$ proves \eqref{eq:cf} for $M_0$.
For each fixed coefficient only finitely many heights occur, so the
infinite fraction is well-defined as a formal power series.
\end{proof}

The shift of every diagonal coefficient by $t$ is also the general
binomial-transform rule proved by Barry~\cite[Proposition 4]{barry-cf}.
His Proposition~22 and its following weighted formula, pp.~17--18,
give the same even-index transform rule for a sequence with a known
Stieltjes fraction; thus that transformation rule is prior work as well.

In particular the Hankel product is independent of $t$. The ordinary
series at $t=1$ has radius zero: $a_{2m}(1)\geq m!$.
Thus \eqref{eq:cf} asserts formal coefficient identities, not convergence
of a power series near zero. Also $a_n'(t)=na_{n-1}(t)$ follows by
differentiating \eqref{eq:egf}, so this is an Appell family.

\section{The asymptotic formula and its weighted extension}
\begin{corollary}\label{cor:asymptotic}
For each fixed real $t>0$, as $n\to\infty$ through all integers,
\begin{equation}\label{eq:asymptotic}
 a_n(t)\sim\frac{\sqrt{\pi n}}2
       \left(\frac n{2e}\right)^{n/2}
       \exp\left(t\sqrt{2n}-\frac{t^2}{2}\right).
\end{equation}
At $t=1$ this is Rousseau's 14 March 2006 formula in A084261.
Fang, Hamaker and Troyka reproduce its even-index specialization in
Table~2, arXiv v3, p.~17, explicitly attributing it to OEIS~\cite{fang}.
\end{corollary}
\begin{proof}
Split \eqref{eq:moment} into positive and negative $x$.
The positive part is
$I_n=\int_0^\infty x e^{\phi_n(x)}\,dx$, with
$\phi_n(x)=n\log(t+x)-x^2$.
Its unique maximum is at
$r=(\sqrt{t^2+2n}-t)/2$, satisfying $n=2r(r+t)$.
For fixed $y$, Taylor expansion at $r$ gives
\[
 \phi_n(r+y)-\phi_n(r)\longrightarrow-2y^2.
\]
Since $\phi_n''(x)=-n/(t+x)^2-2\leq-2$ on $x\geq0$,
the difference is at most $-y^2$. For $r\geq1$, the integrand after
putting $x=r+y$ and dividing by $r e^{\phi_n(r)}$ is bounded by
$(1+|y|)e^{-y^2}$ on $y\geq-r$, and is zero otherwise.
Dominated convergence therefore gives
\begin{equation}\label{eq:laplace}
 I_n\sim r e^{\phi_n(r)}\sqrt{\pi/2}.
\end{equation}

The absolute value of the negative part is bounded by
$J_n=\int_0^\infty x|t-x|^ne^{-x^2}\,dx$.
The interval $[0,t]$ contributes at most $t^{n+2}/2=o(I_n)$.
For large $n$ and $t\leq x\leq2r$,
\[
 \left(\frac{x-t}{x+t}\right)^n
 \leq\exp\left(-\frac{2nt}{2r+t}\right)\longrightarrow0.
\]
Thus that interval contributes $o(I_n)$ as well.
For $x\geq2r$, $|x-t|\leq x+t$ and the concavity bound gives
\[
 \int_{2r}^\infty x e^{\phi_n(x)}\,dx
 \leq e^{\phi_n(r)}\int_r^\infty(r+y)e^{-y^2}\,dy
 =o(r e^{\phi_n(r)}).
\]
Consequently $J_n=o(I_n)$, for either parity of $n$.
Finally, writing $q=\sqrt{n/2}$, expansion of the displayed value of
$r$ gives
\[
 r=q-t/2+t^2/(8q)+O(q^{-3}),\qquad
 \phi_n(r)=\frac n2\log\frac n2-\frac n2
              +t\sqrt{2n}-\frac{t^2}{2}+O(n^{-1/2}).
\]
Substitution in \eqref{eq:laplace} proves \eqref{eq:asymptotic}.
\end{proof}

The condition $t>0$ is essential for this one-sided leading term.
At $t=0$, $a_{2m}(0)=m!$ and $a_{2m+1}(0)=0$; these exact values
are not asserted to satisfy \eqref{eq:asymptotic}.
For negative $t$, the polynomial identity
$a_n(-t)=(-1)^na_n(t)$ gives the corresponding signed formula.

\begin{corollary}\label{cor:clans}
Let $c_n$ count the $2143$-I-avoiding involutions on $n$ elements with
each fixed point independently assigned either of two signs, summing over
all sign imbalances. Then $c_n=a_n(2)$, $c_0=1$, $c_1=2$, $c_2=5$,
and for $n\geq3$,
\[
 2c_n-8c_{n-1}+(8-n)c_{n-2}+2(n-2)c_{n-3}=0.
\]
Moreover $c_n\sim(\sqrt{\pi n}/2)(n/(2e))^{n/2}
\exp(2\sqrt{2n}-2)$.
\end{corollary}
\begin{proof}
An involution with $j$ fixed points has $2^j$ sign assignments.
Use the fixed-point interpretation of $a_n(t)$ and substitute $t=2$
in Corollaries~\ref{cor:weighted} and \ref{cor:asymptotic}.
This specializes the clan interpretation noted by Fang, Hamaker and
Troyka~\cite[Section 5.1]{fang}; no individual sign imbalance is fixed.
\end{proof}

\section{Coverage and verification}
The FORMULA field of A084261 contains several statements. Their relation
to this note is recorded below; distinct formulas are not conflated with
the recurrence conjecture.

\begingroup\small
\begin{longtable}{@{}p{0.43\linewidth}p{0.50\linewidth}@{}}
\toprule
Source item & Coverage \\
\midrule
Inclusive even-index factorial sum & Definition \eqref{eq:sum}; prior
enumeration proved by Fang--Hamaker--Troyka (2020, 2022). \\
Parity-filtered sum & Same sum, with odd summands zero; the parity
sequence $b$ makes the convention explicit. \\
Mathar's 2012 recurrence & Proved for all $n\geq3$ in Theorem~\ref{thm:main}. \\
Jovovic's 25 September 2003 EGF & Reproved in
Corollary~\ref{cor:weighted}; earlier formula credited. \\
Hanna's 17 January 2006 fraction & Proved in
Corollary~\ref{cor:hankel}; earlier formula credited;
prior transform proof Barry (2009), Proposition~22. \\
Rousseau's 14 March 2006 asymptotic & Proved for all indices in
Corollary~\ref{cor:asymptotic}, extended to $t>0$; earlier formula credited. \\
Barry's 11 April 2010 row-sum and Hankel comments & Proved in
Corollaries~\ref{cor:triangle} and \ref{cor:hankel}; earlier entries credited. \\
Fang--Hamaker--Troyka's weighted EGF & Reproved; prior proof
Fang--Hamaker--Troyka (2020; 2022), Proposition~4.13 in v3. \\
A161556 row/column formulas and bivariate fraction & Proved in
Corollaries~\ref{cor:triangle} and \ref{cor:hankel}. \\
A161556 displayed production matrix; links A155856, A156367 &
Not covered. \\
A137704 product and aerated-factorial Hankel description & Proved in
Corollary~\ref{cor:hankel}, including $t=0$. \\
A137704: aerated A003319; Kotesovec's 2024 determinant asymptotic &
Not covered. \\
A002896 lattice-walk formulas (linked beside Rousseau's asymptotic) &
Not covered; this is a different sequence. \\
Fang--Hamaker--Troyka's other pattern classes and open directions &
Not covered. \\
Fang--Hamaker--Troyka, Section~5.1: two signs on fixed points &
Prior interpretation cited; recurrence and asymptotic consequences
proved in Corollary~\ref{cor:clans}. \\
Weighted recurrence, moment positivity and Appell identity &
Proved consequences; not posed as conjectures in A084261. \\
\bottomrule
\end{longtable}
\endgroup

The word ``Proved'' in the table applies to formulas actually stated
in the indicated entries. The weighted extensions are additional
consequences, not refutations or answers to unposed questions.

The Lean declaration
\href{https://github.com/the-omega-institute/trureturing/blob/b0f7c01b01511ea00428c1b6edad5dd2b372eaf8/D5/S1/Recurrence/MatharA084261.lean}
{\texttt{D5.S1.Recurrence.MatharA084261.result}}
proves exactly \eqref{eq:sum}--\eqref{eq:rec} for all natural $n\geq3$,
with integer arithmetic in the recurrence. The generating-function identities
and the polynomial recurrence, moment consequences and asymptotics have
the written proofs above; those claims are distinct from this declaration.
The accompanying Python checker compares the defining sum with the recurrence
for $0\leq n\leq880$, and tests \eqref{eq:weightedrec} for
$t\in\{-2,-1,0,1,2\}$ and $3\leq n\leq120$ using exact integers.
It also checks row polynomials, the Appell identity, finite continued-fraction
coefficients and Hankel determinants. These finite checks complement
the universal proofs; the asymptotic claim rests on its written argument.
The source package contains this text, the PDF, metadata and the exact
checker with its expected output and rebuild instructions.

\begin{thebibliography}{9}
\bibitem{oeis}
The OEIS Foundation Inc.,
\emph{A084261: A binomial transform of factorial numbers}.
Entry by Paul Barry, 26 May 2003; exponential generating function by
Vladeta Jovovic, 25 September 2003; recurrence conjecture by
R.\ J.\ Mathar, 30 November 2012.
\url{https://oeis.org/A084261}. Accessed 8 November 2026.

\bibitem{fang}
Jonathan J. Fang, Zachary Hamaker and Justin M. Troyka,
\emph{On pattern avoidance in matchings and involutions},
\emph{Electronic Journal of Combinatorics} \textbf{29}(1) (2022), P1.39,
published 25 February 2022. \url{https://doi.org/10.37236/10155}.
First version: arXiv:2009.00079v1, 31 August 2020, Theorem~1.7 and
Section~4.3. Version~3, 12 January 2022, Section~4.3,
Proposition~4.13 and following example, p.~16; Table~2, p.~17.
\url{https://arxiv.org/abs/2009.00079v3}.

\bibitem{oeis-triangle}
The OEIS Foundation Inc., \emph{A161556: Exponential Riordan array
$[1+(\sqrt\pi/2)x e^{x^2/4}\operatorname{erf}(x/2),x]$}.
Entry by Paul Barry, 13 June 2009.
\url{https://oeis.org/A161556}. Accessed 8 November 2026.

\bibitem{oeis-hankel}
The OEIS Foundation Inc., \emph{A137704: Hankel transform of aerated
factorial numbers}. Entry by Paul Barry, 7 February 2008;
asymptotic formula by Vaclav Kotesovec, 20 January 2024.
\url{https://oeis.org/A137704}. Accessed 8 November 2026.

\bibitem{dlmf}
NIST Digital Library of Mathematical Functions, Chapter~18,
\emph{Orthogonal Polynomials}, Sections~18.3 and 18.5.
\url{https://dlmf.nist.gov/18.3}; \url{https://dlmf.nist.gov/18.5}.
Accessed 8 November 2026.

\bibitem{barry-cf}
Paul Barry, \emph{Continued Fractions and Transformations of Integer
Sequences}, Journal of Integer Sequences \textbf{12} (2009), Article 09.7.6,
Propositions~4 and 22, pp.~5--6 and 17--18. Published 4 November 2009.
\url{https://cs.uwaterloo.ca/journals/JIS/VOL12/Barry3/barry93.pdf}.
\end{thebibliography}

\bigskip\noindent\small
\copyright\ 2026 Haobo Ma and Wenlin Zhang.
Licensed under Creative Commons Attribution 4.0 International (CC BY 4.0).
\url{https://creativecommons.org/licenses/by/4.0/}
\end{document}
